% SOURCE_STATEMENT: "For the adaptive-resolution interpolation formula in molecular dynamics, why is it not naturally formulated using potentials, and can the force-based interpolation formula be formally derived from first principles of statistical physics?"
% Preparation and final self-audit: CONTRIBUTING.md prompts C and D; source/audit artifacts retained in the local physics_statements_2026_10_11 batch.
% !TeX program = lualatex
\documentclass[11pt]{article}
\usepackage[margin=25mm]{geometry}
\usepackage{fontspec}
\setmainfont{Latin Modern Roman}
\usepackage{amsmath,amssymb,mathtools}
\usepackage{unicode-math}
\setmathfont{Latin Modern Math}
\usepackage{xurl,hyperref}
\hypersetup{colorlinks=true,urlcolor=blue}
\setcounter{secnumdepth}{0}
\setlength{\parindent}{0pt}
\setlength{\parskip}{5pt}
\setlength{\emergencystretch}{3em}
\newcommand{\sref}[2]{\hyperlink{source:#1:#2}{[S#2]}}
\newcommand{\cataloguescope}[1]{\paragraph{Recorded target.}#1}
\newcommand{\R}{\mathbb R}
\newcommand{\C}{\mathbb C}
\newcommand{\N}{\mathbb N}
\newcommand{\Z}{\mathbb Z}
\newcommand{\Q}{\mathbb Q}
\newcommand{\F}{\mathbb F}
\newcommand{\eps}{\varepsilon}
\begin{document}
\section{First-Principles Derivation of Adaptive-Resolution Molecular Dynamics}\label{statement}

\subsection{Definitions and mathematical statement}

\paragraph{Shared notation.}$\mathbb N=\{1,2,\ldots\}$; the standard number fields have their usual meanings. All additional conventions are stated locally.


\paragraph{Definitions and conventions.} Adaptive-resolution molecular dynamics couples an atomistic region and a coarse-grained region using a switching function $w\in[0,1]$. For molecular centers $X_\alpha,X_\beta$, the report's force interpolation has atomistic weight $w(X_\alpha)w(X_\beta)$ and coarse-grained weight $1-w(X_\alpha)w(X_\beta)$. A potential interpolation generally yields additional derivatives of the spatial weights, hence drift forces. A first-principles derivation would obtain the prescribed interpolation and statistical ensemble from microscopic statistical mechanics rather than choosing them ad hoc.

\paragraph{Source scope and qualifications.} The report supplies an argument against a particular potential-interpolation scheme and leaves a first-principles force derivation open. It does not establish that every conceivable energy-based coarse-graining is impossible; that stronger conclusion is not part of the target.

\paragraph{Catalogue mathematical target.}
\paragraph{Catalogue wording retained for traceability.}
For the adaptive-resolution interpolation formula in molecular dynamics, why is it not naturally formulated using potentials, and can the force-based interpolation formula be formally derived from first principles of statistical physics?

\paragraph{Saved literature context.} The supplied catalogue's assessment dated 2026-08-21 records: The potential-based coupling is ruled out by an overdetermined drift system in the report, while a first-principles derivation of the force interpolation formula is explicitly left as a challenge. The assessment is retained as dated provenance; it is not a new comprehensive literature or open-status certification.
\subsection{Short English statement}
Explain why the adaptive-resolution coupling is force based and derive that force interpolation from statistical mechanics.
\subsection{Sources}
\begin{itemize}
\item[\hypertarget{source:statement:1}{S1}] ULAM.AI and the UnsolvedMath Contributors, UnsolvedMath, catalogue entry. Catalogue curation under CC BY 4.0; definitions and English exposition added locally.
\url{https://huggingface.co/datasets/ulamai/UnsolvedMath}.
\item[\hypertarget{source:statement:2}{S2}] Atomistic Models of Materials: Mathematical Challenges (2008), Oberwolfach Reports, DOI 10.4171/owr/2008/21, pp. 1079--1156; official report source and mathematical context.
\url{https://doi.org/10.4171/owr/2008/21}.
\item[\hypertarget{source:statement:3}{S3}] Official EMS Press PDF of the same Oberwolfach report; PDF pages 28--30.
\url{https://ems.press/content/serial-article-files/46165}.
\item[\hypertarget{source:statement:4}{S4}] Delle Site 2007; reference in the saved dated catalogue literature review.
\url{https://doi.org/10.1103/PhysRevE.76.047701}.
\end{itemize}
\label{end:statement}
\end{document}
